% =========================================================
% CHUONG: MANG NO-RON NHAN TAO (NEURAL NETWORK)
% =========================================================

\chapter{Mạng nơ-ron nhân tạo \\ (Neural Network)}
\label{chap:neural-network}

Mạng nơ-ron nhân tạo là họ mô hình đứng sau hầu hết thành tựu nổi bật của trí tuệ nhân tạo hiện đại: nhận dạng ảnh, dịch máy, nhận dạng giọng nói và các mô hình ngôn ngữ lớn. Chương cuối này trình bày nền tảng của mạng nơ-ron: từ perceptron đơn lẻ đến mạng nhiều lớp (Multilayer Perceptron -- MLP), thuật toán lan truyền ngược, các kỹ thuật huấn luyện, và cách cài đặt bằng scikit-learn cùng một thư viện học sâu.

\noindent Sau khi học xong chương này, người học có khả năng:
\begin{itemize}
    \item Trình bày lịch sử phát triển của mạng nơ-ron và lý do của các ``mùa đông AI''.
    \item Tính tay đầu ra của perceptron và giải thích vì sao perceptron không học được XOR.
    \item Tính tay một bước lan truyền xuôi và lan truyền ngược trên mạng nhỏ.
    \item Huấn luyện \texttt{MLPClassifier}/\texttt{MLPRegressor} và một mạng Keras đơn giản, lựa chọn hàm kích hoạt, hàm mất mát, bộ tối ưu và kỹ thuật điều chuẩn.
\end{itemize}

% ---------------------------------------------------------
\section{Lịch sử phát triển}
% ---------------------------------------------------------

\begin{table}[H]
\centering
\begin{tabular}{|c|p{11.5cm}|}
\hline
\textbf{Năm} & \textbf{Sự kiện} \\
\hline
1943 & McCulloch và Pitts đề xuất mô hình toán học đầu tiên của nơ-ron sinh học: nơ-ron ngưỡng với đầu vào nhị phân. \\ \hline
1949 & Donald Hebb nêu quy tắc học ``các nơ-ron cùng kích hoạt thì liên kết mạnh lên''. \\ \hline
1958 & Frank Rosenblatt công bố \textbf{Perceptron} --- mô hình nơ-ron đầu tiên có thuật toán học trọng số từ dữ liệu, cài đặt trên phần cứng Mark I. \\ \hline
1969 & Minsky và Papert (sách \textit{Perceptrons}) chứng minh perceptron một lớp không biểu diễn được hàm XOR $\Rightarrow$ ``mùa đông AI'' thứ nhất. \\ \hline
1986 & Rumelhart, Hinton và Williams phổ biến thuật toán \textbf{lan truyền ngược} (backpropagation) cho mạng nhiều lớp (ý tưởng đã có từ Linnainmaa 1970, Werbos 1974). \\ \hline
1989 & Cybenko, Hornik chứng minh \textbf{định lý xấp xỉ vạn năng}. LeCun áp dụng mạng tích chập (CNN) nhận dạng chữ số viết tay. \\ \hline
1997 & Hochreiter và Schmidhuber đề xuất LSTM cho dữ liệu chuỗi. \\ \hline
2006 & Hinton và cộng sự huấn luyện mạng sâu bằng tiền huấn luyện từng lớp --- khởi đầu thuật ngữ ``deep learning''. \\ \hline
2012 & AlexNet (Krizhevsky, Sutskever, Hinton) thắng ImageNet với cách biệt lớn nhờ GPU, ReLU, Dropout. \\ \hline
2017 & Vaswani và cộng sự giới thiệu kiến trúc \textbf{Transformer} (``Attention Is All You Need''), nền tảng của các mô hình ngôn ngữ lớn. \\ \hline
2018 & Bengio, Hinton, LeCun nhận giải Turing cho những đóng góp về học sâu. \\ \hline
\end{tabular}
\caption{Các mốc lịch sử của mạng nơ-ron}
\end{table}

% ---------------------------------------------------------
\section{Perceptron}
% ---------------------------------------------------------

\begin{definition}[Perceptron]
Perceptron nhận vector đầu vào $\mathbf{x} \in \mathbb{R}^d$, tính tổng có trọng số rồi đưa qua hàm bước:
\begin{equation}
    \hat{y} = \operatorname{step}(\mathbf{w}^{\top}\mathbf{x} + b), \qquad \operatorname{step}(z) = \begin{cases} 1 & z \geq 0 \\ 0 & z < 0 \end{cases}
\end{equation}
\end{definition}

\noindent\textbf{Quy tắc học của Rosenblatt:} với mỗi mẫu bị phân lớp sai, cập nhật
\begin{equation}
    \mathbf{w} \leftarrow \mathbf{w} + \eta\,(y - \hat{y})\,\mathbf{x}, \qquad b \leftarrow b + \eta\,(y - \hat{y}).
\end{equation}
\textbf{Định lý hội tụ perceptron} (Novikoff, 1962): nếu dữ liệu \textbf{tách được tuyến tính}, thuật toán dừng sau hữu hạn bước cập nhật.

\hinhchuachen{ch-nn-perceptron.png}{Cấu trúc perceptron: đầu vào, trọng số, tổng có trọng số, hàm kích hoạt và đầu ra.}{fig:perceptron}

\begin{example}[Perceptron cho hàm AND]
Chọn $\mathbf{w} = (1, 1)$, $b = -1{,}5$:
\end{example}

\begin{table}[H]
\centering
\begin{tabular}{|c|c|c|c|c|c|}
\hline
$x_1$ & $x_2$ & $z = x_1 + x_2 - 1{,}5$ & $\hat{y}$ & AND & XOR \\
\hline
0 & 0 & $-1{,}5$ & 0 & 0 & 0 \\ \hline
0 & 1 & $-0{,}5$ & 0 & 0 & 1 \\ \hline
1 & 0 & $-0{,}5$ & 0 & 0 & 1 \\ \hline
1 & 1 & $0{,}5$ & 1 & 1 & 0 \\ \hline
\end{tabular}
\caption{Perceptron biểu diễn hàm AND}
\label{tab:and}
\end{table}

Đường thẳng $x_1 + x_2 = 1{,}5$ tách điểm $(1, 1)$ khỏi ba điểm còn lại. Nhưng với \textbf{XOR}, hai điểm lớp 1 là $(0, 1)$, $(1, 0)$ nằm trên một đường chéo và hai điểm lớp 0 nằm trên đường chéo kia: \textbf{không đường thẳng nào} tách được. Thật vậy, giả sử tồn tại $\mathbf{w}, b$: từ $(0, 0) \mapsto 0$ có $b < 0$; từ $(0, 1), (1, 0) \mapsto 1$ có $w_2 + b \geq 0$, $w_1 + b \geq 0$, cộng lại được $w_1 + w_2 + b \geq -b > 0$, mâu thuẫn với $(1, 1) \mapsto 0$. Lời giải: thêm một \textbf{lớp ẩn} --- XOR $= $ AND(OR, NAND).

\begin{lstlisting}[caption={Cài đặt perceptron từ đầu}]
import numpy as np

class Perceptron:
    def __init__(self, lr=1.0, epochs=10):
        self.lr, self.epochs = lr, epochs

    def fit(self, X, y):
        self.w = np.zeros(X.shape[1]); self.b = 0.0
        for _ in range(self.epochs):
            errors = 0
            for xi, yi in zip(X, y):
                update = self.lr * (yi - self.predict(xi))
                self.w += update * xi
                self.b += update
                errors += int(update != 0)
            if errors == 0:          # hoi tu
                break
        return self

    def predict(self, X):
        return (np.dot(X, self.w) + self.b >= 0).astype(int)

X = np.array([[0, 0], [0, 1], [1, 0], [1, 1]])
p_and = Perceptron().fit(X, np.array([0, 0, 0, 1]))
print(p_and.w, p_and.b, p_and.predict(X))   # hoc dung AND
p_xor = Perceptron(epochs=100).fit(X, np.array([0, 1, 1, 0]))
print(p_xor.predict(X))                     # khong bao gio dung ca 4
\end{lstlisting}

% ---------------------------------------------------------
\section{Mạng nơ-ron nhiều lớp (MLP)}
% ---------------------------------------------------------

\subsection{Kiến trúc}

MLP gồm \textbf{lớp đầu vào}, một hoặc nhiều \textbf{lớp ẩn} và \textbf{lớp đầu ra}; mỗi nơ-ron nối với mọi nơ-ron ở lớp kế tiếp (lớp kết nối đầy đủ -- dense). Với lớp thứ $l$:
\begin{equation}
    \mathbf{z}^{(l)} = \mathbf{W}^{(l)}\mathbf{a}^{(l-1)} + \mathbf{b}^{(l)}, \qquad \mathbf{a}^{(l)} = g^{(l)}\big(\mathbf{z}^{(l)}\big), \qquad \mathbf{a}^{(0)} = \mathbf{x}.
\end{equation}
Hàm kích hoạt $g$ phải \textbf{phi tuyến}: nếu không, nhiều lớp tuyến tính chồng lên nhau vẫn chỉ là một phép biến đổi tuyến tính.

\hinhchuachen{ch-nn-mlp.png}{Mạng MLP với lớp đầu vào, hai lớp ẩn và lớp đầu ra.}{fig:mlp}

\begin{theorem}[Xấp xỉ vạn năng -- Cybenko 1989, Hornik 1991]
Mạng một lớp ẩn với đủ nhiều nơ-ron và hàm kích hoạt phi tuyến (không phải đa thức) có thể xấp xỉ mọi hàm liên tục trên tập compact với độ chính xác tùy ý.
\end{theorem}
Định lý chỉ khẳng định \textbf{tồn tại}, không nói cách tìm trọng số hay cần bao nhiêu nơ-ron. Trong thực tế, mạng \textbf{sâu} thường biểu diễn hiệu quả hơn mạng rộng một lớp.

\subsection{Hàm kích hoạt}

\begin{table}[H]
\centering
\begin{tabular}{|l|c|c|p{4.8cm}|}
\hline
\textbf{Hàm} & \textbf{Công thức} & \textbf{Đạo hàm} & \textbf{Ghi chú} \\
\hline
Sigmoid & $\sigma(z) = \frac{1}{1 + e^{-z}}$ & $\sigma(1 - \sigma)$ & Đầu ra $(0, 1)$; bão hòa $\Rightarrow$ triệt tiêu gradient \\[4pt] \hline
Tanh & $\tanh z$ & $1 - \tanh^2 z$ & Đầu ra $(-1, 1)$, tâm 0 \\[4pt] \hline
ReLU & $\max(0, z)$ & $\mathbb{1}[z > 0]$ & Mặc định cho lớp ẩn (Nair \& Hinton, 2010); nhanh \\[4pt] \hline
Leaky ReLU & $\max(0{,}01z, z)$ & $1$ hoặc $0{,}01$ & Tránh nơ-ron ``chết'' \\[4pt] \hline
Softmax & $\frac{e^{z_k}}{\sum_j e^{z_j}}$ & --- & Lớp đầu ra phân lớp nhiều lớp \\[4pt] \hline
\end{tabular}
\caption{Các hàm kích hoạt phổ biến}
\end{table}

\hinhchuachen{ch-nn-activations.png}{Đồ thị các hàm kích hoạt sigmoid, tanh, ReLU và Leaky ReLU.}{fig:activations}

\subsection{Hàm mất mát và lớp đầu ra}

\begin{table}[H]
\centering
\begin{tabular}{|l|l|l|}
\hline
\textbf{Bài toán} & \textbf{Lớp đầu ra} & \textbf{Hàm mất mát} \\
\hline
Hồi quy & 1 nơ-ron, tuyến tính & MSE \\ \hline
Phân lớp nhị phân & 1 nơ-ron, sigmoid & Binary cross-entropy \\ \hline
Phân lớp $K$ lớp & $K$ nơ-ron, softmax & Categorical cross-entropy $-\sum_k y_k \log \hat{y}_k$ \\ \hline
\end{tabular}
\caption{Ghép cặp lớp đầu ra và hàm mất mát}
\end{table}

% ---------------------------------------------------------
\section{Lan truyền ngược (Backpropagation)}
% ---------------------------------------------------------

Huấn luyện mạng là cực tiểu hàm mất mát $L$ theo mọi trọng số bằng gradient descent (Chương~\ref{chap:regression}). Lan truyền ngược là cách tính \textbf{mọi} đạo hàm riêng $\partial L / \partial w$ hiệu quả bằng \textbf{quy tắc dây chuyền}, đi từ lớp đầu ra về lớp đầu vào và tái sử dụng kết quả trung gian. Chi phí chỉ cỡ một lần lan truyền xuôi. Định nghĩa sai số tại lớp $l$ là $\boldsymbol{\delta}^{(l)} = \partial L / \partial \mathbf{z}^{(l)}$:
\begin{align}
    \boldsymbol{\delta}^{(L)} &= \nabla_{\mathbf{a}^{(L)}} L \odot g'\big(\mathbf{z}^{(L)}\big), \\
    \boldsymbol{\delta}^{(l)} &= \big(\mathbf{W}^{(l+1)}\big)^{\top}\boldsymbol{\delta}^{(l+1)} \odot g'\big(\mathbf{z}^{(l)}\big), \\
    \frac{\partial L}{\partial \mathbf{W}^{(l)}} &= \boldsymbol{\delta}^{(l)}\big(\mathbf{a}^{(l-1)}\big)^{\top}, \qquad
    \frac{\partial L}{\partial \mathbf{b}^{(l)}} = \boldsymbol{\delta}^{(l)} .
\end{align}

\begin{example}[Tính tay một bước lan truyền xuôi và ngược]
Mạng 2--2--1, mọi nơ-ron dùng sigmoid, mất mát $L = \frac{1}{2}(t - o)^2$. Đầu vào $\mathbf{x} = (1;\ 0{,}5)$, nhãn $t = 1$,
\[
    \mathbf{W}^{(1)} = \begin{pmatrix} 0{,}1 & 0{,}2 \\ 0{,}3 & 0{,}4 \end{pmatrix}, \quad \mathbf{b}^{(1)} = \begin{pmatrix} 0{,}1 \\ 0{,}1 \end{pmatrix}, \quad
    \mathbf{w}^{(2)} = (0{,}5;\ 0{,}6), \quad b^{(2)} = 0{,}1, \quad \eta = 0{,}5 .
\]
\end{example}

\noindent\textbf{Lan truyền xuôi:}
\begin{align*}
    z_1 &= 0{,}1 \cdot 1 + 0{,}2 \cdot 0{,}5 + 0{,}1 = 0{,}3, & h_1 &= \sigma(0{,}3) = 0{,}5744, \\
    z_2 &= 0{,}3 \cdot 1 + 0{,}4 \cdot 0{,}5 + 0{,}1 = 0{,}6, & h_2 &= \sigma(0{,}6) = 0{,}6457, \\
    z_o &= 0{,}5 \cdot 0{,}5744 + 0{,}6 \cdot 0{,}6457 + 0{,}1 = 0{,}7746, & o &= \sigma(0{,}7746) = 0{,}6845,
\end{align*}
mất mát $L = \frac{1}{2}(1 - 0{,}6845)^2 = 0{,}0498$.

\noindent\textbf{Lan truyền ngược:}
\begin{align*}
    \delta_o &= (o - t)\,o(1 - o) = (-0{,}3155)(0{,}6845)(0{,}3155) = -0{,}0681, \\
    \frac{\partial L}{\partial \mathbf{w}^{(2)}} &= \delta_o\,\mathbf{h} = (-0{,}0391;\ -0{,}0440), \\
    \delta_{h_1} &= \delta_o \cdot 0{,}5 \cdot h_1(1 - h_1) = -0{,}0083, \qquad
    \delta_{h_2} = \delta_o \cdot 0{,}6 \cdot h_2(1 - h_2) = -0{,}0094, \\
    \frac{\partial L}{\partial \mathbf{W}^{(1)}} &= \boldsymbol{\delta}_h\,\mathbf{x}^{\top} = \begin{pmatrix} -0{,}0083 & -0{,}0042 \\ -0{,}0094 & -0{,}0047 \end{pmatrix}.
\end{align*}

\noindent\textbf{Cập nhật:} $\mathbf{w}^{(2)} \leftarrow (0{,}5;\ 0{,}6) - 0{,}5 \times (-0{,}0391;\ -0{,}0440) = (0{,}5196;\ 0{,}6220)$. Gradient âm nên trọng số tăng, đẩy $o$ lên gần nhãn $t = 1$. Chú ý gradient ở lớp ẩn nhỏ hơn lớp ra khoảng 8 lần, vì mỗi lần nhân với $\sigma' \leq 0{,}25$ --- đây chính là hiện tượng \textbf{triệt tiêu gradient} khiến mạng sâu dùng sigmoid khó huấn luyện, và là lý do ReLU được ưa chuộng.

\begin{lstlisting}[caption={Kiểm tra ví dụ lan truyền ngược bằng NumPy}]
sigmoid = lambda z: 1 / (1 + np.exp(-z))
x = np.array([1.0, 0.5]); t = 1.0
W1 = np.array([[0.1, 0.2], [0.3, 0.4]]); b1 = np.array([0.1, 0.1])
w2 = np.array([0.5, 0.6]); b2 = 0.1; lr = 0.5

# Lan truyen xuoi
h = sigmoid(W1 @ x + b1)            # [0.5744 0.6457]
o = sigmoid(w2 @ h + b2)            # 0.6845
loss = 0.5 * (t - o) ** 2           # 0.0498

# Lan truyen nguoc
d_o = (o - t) * o * (1 - o)         # -0.0681
d_h = d_o * w2 * h * (1 - h)        # [-0.0083 -0.0094]
w2 -= lr * d_o * h;  b2 -= lr * d_o
W1 -= lr * np.outer(d_h, x);  b1 -= lr * d_h
print(w2)                           # [0.5196 0.622 ]
\end{lstlisting}

% ---------------------------------------------------------
\section{Huấn luyện mạng nơ-ron trong thực tế}
% ---------------------------------------------------------

\subsection{Bộ tối ưu}
\begin{itemize}
    \item \textbf{SGD theo lô nhỏ} (mini-batch): mỗi bước dùng 32--256 mẫu; một \textbf{epoch} là một lượt duyệt hết dữ liệu.
    \item \textbf{Momentum}: cộng dồn hướng các bước trước, giúp vượt vùng phẳng và giảm dao động.
    \item \textbf{Adam} (Kingma \& Ba, 2015): tốc độ học thích nghi cho từng tham số; lựa chọn mặc định tốt (\texttt{solver="adam"} trong scikit-learn).
\end{itemize}
Tốc độ học $\eta$ là siêu tham số quan trọng nhất: quá lớn thì phân kỳ, quá nhỏ thì hội tụ chậm.

\subsection{Khởi tạo và chuẩn hóa}
\begin{itemize}
    \item Không khởi tạo mọi trọng số bằng 0 (các nơ-ron sẽ học giống hệt nhau). Dùng khởi tạo Xavier/Glorot (2010) cho sigmoid/tanh, He (2015) cho ReLU.
    \item \textbf{Chuẩn hóa đầu vào} (\texttt{StandardScaler}) là bắt buộc --- mạng nơ-ron rất nhạy với thang đo.
    \item \textbf{Batch Normalization} (Ioffe \& Szegedy, 2015) chuẩn hóa đầu ra các lớp ẩn, giúp huấn luyện nhanh và ổn định.
\end{itemize}

\subsection{Chống quá khớp}
\begin{itemize}
    \item \textbf{Điều chuẩn L2} (weight decay): tham số \texttt{alpha} trong \texttt{MLPClassifier}.
    \item \textbf{Dropout} (Srivastava và cộng sự, 2014): mỗi bước huấn luyện tắt ngẫu nhiên một tỉ lệ nơ-ron.
    \item \textbf{Dừng sớm} (early stopping): theo dõi lỗi validation, dừng khi không còn cải thiện.
    \item Tăng cường dữ liệu (data augmentation) với ảnh, âm thanh.
\end{itemize}

\hinhchuachen{ch-nn-loss-curve.png}{Đường cong mất mát huấn luyện và kiểm định theo epoch; điểm dừng sớm tại nơi lỗi kiểm định nhỏ nhất.}{fig:loss-curve}

% ---------------------------------------------------------
\section{MLP với scikit-learn}
% ---------------------------------------------------------

\begin{lstlisting}[caption={MLPClassifier giải bài toán XOR}]
from sklearn.neural_network import MLPClassifier

X = np.array([[0, 0], [0, 1], [1, 0], [1, 1]]); y = np.array([0, 1, 1, 0])
mlp = MLPClassifier(hidden_layer_sizes=(4,), activation="tanh",
                    solver="lbfgs", random_state=1, max_iter=1000)
mlp.fit(X, y)
print(mlp.predict(X))     # [0 1 1 0] -- mot lop an du de hoc XOR
\end{lstlisting}

\begin{lstlisting}[caption={Phân lớp chữ số viết tay với MLPClassifier}]
from sklearn.datasets import load_digits
from sklearn.model_selection import train_test_split, GridSearchCV
from sklearn.pipeline import make_pipeline
from sklearn.preprocessing import StandardScaler
from sklearn.metrics import classification_report
import matplotlib.pyplot as plt

X, y = load_digits(return_X_y=True)
X_tr, X_te, y_tr, y_te = train_test_split(X, y, test_size=0.2,
                                          stratify=y, random_state=42)
pipe = make_pipeline(
    StandardScaler(),
    MLPClassifier(hidden_layer_sizes=(128, 64), activation="relu",
                  solver="adam", alpha=1e-4, learning_rate_init=1e-3,
                  batch_size=64, max_iter=300, early_stopping=True,
                  random_state=42))
pipe.fit(X_tr, y_tr)
print(classification_report(y_te, pipe.predict(X_te), digits=4))

mlp = pipe[-1]
plt.plot(mlp.loss_curve_, label="train loss")
plt.plot(mlp.validation_scores_, label="validation accuracy")
plt.legend(); plt.xlabel("epoch"); plt.show()

# Tim sieu tham so (Chuong Validation)
grid = GridSearchCV(pipe, {"mlpclassifier__hidden_layer_sizes": [(64,), (128, 64)],
                           "mlpclassifier__alpha": [1e-4, 1e-2, 1]}, cv=3, n_jobs=-1)
grid.fit(X_tr, y_tr); print(grid.best_params_, grid.score(X_te, y_te))
\end{lstlisting}

\begin{lstlisting}[caption={MLPRegressor dự đoán giá nhà Ames}]
from sklearn.neural_network import MLPRegressor
from sklearn.compose import TransformedTargetRegressor

# preprocess: ColumnTransformer da xay dung o Chuong Tien xu ly du lieu
reg = make_pipeline(preprocess,
                    MLPRegressor(hidden_layer_sizes=(64, 32), alpha=1e-3,
                                 early_stopping=True, max_iter=1000, random_state=0))
model = TransformedTargetRegressor(regressor=reg, func=np.log1p, inverse_func=np.expm1)
model.fit(X_train, y_train)
\end{lstlisting}

\noindent\textit{Giới hạn:} \texttt{MLPClassifier} chỉ chạy trên CPU và chỉ hỗ trợ lớp dense --- phù hợp dữ liệu dạng bảng cỡ vừa và mục đích học tập. Với ảnh, văn bản, âm thanh và mạng sâu, cần các thư viện học sâu chuyên dụng.

% ---------------------------------------------------------
\section{Giới thiệu thư viện học sâu}
% ---------------------------------------------------------

\textbf{TensorFlow/Keras} (Google, 2015) và \textbf{PyTorch} (Meta, 2016) là hai thư viện phổ biến nhất. Cả hai cung cấp tensor chạy trên GPU, \textbf{tự động tính đạo hàm} (automatic differentiation --- tổng quát hóa của lan truyền ngược) và các lớp dựng sẵn: Dense, Conv2D (CNN cho ảnh), LSTM/GRU (dữ liệu chuỗi), Attention/Transformer (văn bản).

\begin{lstlisting}[caption={Cùng bài toán digits với Keras}]
import keras
from keras import layers

scaler = StandardScaler().fit(X_tr)
model = keras.Sequential([
    keras.Input(shape=(64,)),
    layers.Dense(128, activation="relu"),
    layers.Dropout(0.2),
    layers.Dense(64, activation="relu"),
    layers.Dense(10, activation="softmax"),
])
model.compile(optimizer=keras.optimizers.Adam(1e-3),
              loss="sparse_categorical_crossentropy", metrics=["accuracy"])
model.summary()
stop = keras.callbacks.EarlyStopping(patience=5, restore_best_weights=True)
history = model.fit(scaler.transform(X_tr), y_tr, epochs=50, batch_size=64,
                    validation_split=0.1, callbacks=[stop])
print(model.evaluate(scaler.transform(X_te), y_te))
\end{lstlisting}

\begin{lstlisting}[caption={Vòng lặp huấn luyện tối giản với PyTorch}]
import torch
from torch import nn

net = nn.Sequential(nn.Linear(64, 128), nn.ReLU(), nn.Linear(128, 10))
opt = torch.optim.Adam(net.parameters(), lr=1e-3)
loss_fn = nn.CrossEntropyLoss()
Xt = torch.tensor(scaler.transform(X_tr), dtype=torch.float32)
yt = torch.tensor(y_tr)
for epoch in range(100):
    opt.zero_grad()
    loss = loss_fn(net(Xt), yt)    # lan truyen xuoi
    loss.backward()                # lan truyen nguoc (autograd)
    opt.step()                     # cap nhat trong so
\end{lstlisting}

\begin{table}[H]
\centering
\begin{tabular}{|p{3.2cm}|p{4.8cm}|p{4.8cm}|}
\hline
\textbf{Loại dữ liệu} & \textbf{Kiến trúc phù hợp} & \textbf{Ví dụ ứng dụng} \\
\hline
Dạng bảng & MLP; thường thua Gradient Boosting (Chương~\ref{chap:ensemble}) & Dự đoán giá, rủi ro tín dụng \\ \hline
Ảnh & CNN, Vision Transformer & Phân loại ảnh, phát hiện đối tượng \\ \hline
Chuỗi, văn bản & RNN/LSTM, Transformer & Dịch máy, phân tích cảm xúc, LLM \\ \hline
\end{tabular}
\caption{Lựa chọn kiến trúc theo loại dữ liệu}
\end{table}

% ---------------------------------------------------------
\section{Tóm tắt chương}
% ---------------------------------------------------------
\begin{itemize}
    \item Perceptron là bộ phân lớp tuyến tính, không học được XOR; thêm lớp ẩn với hàm kích hoạt phi tuyến cho MLP có khả năng xấp xỉ vạn năng.
    \item Lan truyền xuôi tính đầu ra; lan truyền ngược dùng quy tắc dây chuyền tính gradient cho mọi trọng số; cập nhật bằng SGD/Adam.
    \item ReLU cho lớp ẩn; sigmoid/softmax cho lớp ra phân lớp; ghép đúng với cross-entropy.
    \item Chuẩn hóa đầu vào, khởi tạo hợp lý, điều chuẩn L2, Dropout và dừng sớm là các kỹ thuật thiết yếu.
    \item scikit-learn cho MLP cơ bản; Keras/PyTorch cho học sâu trên GPU.
\end{itemize}

% ---------------------------------------------------------
\section{Bài tập}
% ---------------------------------------------------------
\begin{enumerate}
    \item Tìm trọng số perceptron cho hàm OR và NAND. Ghép chúng với perceptron AND (Bảng~\ref{tab:and}) thành mạng hai lớp biểu diễn XOR; kiểm tra trên cả 4 đầu vào.
    \item Chạy lại ví dụ lan truyền ngược thêm 100 bước. Vẽ mất mát theo số bước. Lặp lại với hàm kích hoạt ReLU ở lớp ẩn.
    \item Chứng minh $\sigma'(z) = \sigma(z)(1 - \sigma(z))$ và $\sigma'(z) \leq 0{,}25$. Ước lượng độ lớn gradient ở lớp đầu tiên của mạng 10 lớp sigmoid.
    \item Với softmax và cross-entropy, chứng minh $\partial L / \partial z_k = \hat{y}_k - y_k$.
    \item Huấn luyện \texttt{MLPClassifier} trên tập \texttt{digits} có và không có \texttt{StandardScaler}. So sánh độ chính xác và số vòng lặp.
    \item So sánh \texttt{MLPRegressor}, Ridge và Gradient Boosting trên Ames Housing bằng RMSE kiểm định chéo 5-fold.
\end{enumerate}

% ---------------------------------------------------------
\section*{Tài liệu tham khảo của chương}
\addcontentsline{toc}{section}{Tài liệu tham khảo của chương}
% ---------------------------------------------------------
\begin{enumerate}
    \item Cybenko, G. (1989). Approximation by superpositions of a sigmoidal function. \textit{Mathematics of Control, Signals and Systems}, 2(4), 303--314.
    \item Glorot, X., \& Bengio, Y. (2010). Understanding the difficulty of training deep feedforward neural networks. \textit{Proceedings of AISTATS 2010}, 249--256.
    \item Goodfellow, I., Bengio, Y., \& Courville, A. (2016). \textit{Deep Learning}. MIT Press.
    \item Hornik, K. (1991). Approximation capabilities of multilayer feedforward networks. \textit{Neural Networks}, 4(2), 251--257.
    \item Kingma, D. P., \& Ba, J. (2015). Adam: A Method for Stochastic Optimization. \textit{ICLR 2015}.
    \item LeCun, Y., Bengio, Y., \& Hinton, G. (2015). Deep learning. \textit{Nature}, 521, 436--444.
    \item McCulloch, W. S., \& Pitts, W. (1943). A logical calculus of the ideas immanent in nervous activity. \textit{Bulletin of Mathematical Biophysics}, 5, 115--133.
    \item Minsky, M., \& Papert, S. (1969). \textit{Perceptrons}. MIT Press.
    \item Rosenblatt, F. (1958). The perceptron: A probabilistic model for information storage and organization in the brain. \textit{Psychological Review}, 65(6), 386--408.
    \item Rumelhart, D. E., Hinton, G. E., \& Williams, R. J. (1986). Learning representations by back-propagating errors. \textit{Nature}, 323, 533--536.
    \item Srivastava, N., Hinton, G., Krizhevsky, A., Sutskever, I., \& Salakhutdinov, R. (2014). Dropout: A Simple Way to Prevent Neural Networks from Overfitting. \textit{Journal of Machine Learning Research}, 15, 1929--1958.
    \item Vaswani, A., et al. (2017). Attention Is All You Need. \textit{Advances in Neural Information Processing Systems 30}.
\end{enumerate}
